\begin{examplebox}
	\textbf{\textcolor{CExample}{例 1（基础）}}
	\textbf{\textcolor{CExample}{题目：}}
	正方形 $ABCD$ 边长为 $4$，$E$ 在 $BC$ 上，$BE=1$，$\angle EAF=45^\circ$，求 $DF$ 和 $EF$。
	\textbf{\textcolor{CExample}{解题步骤：}}
	\begin{description}[style=nextline, leftmargin=2.8em, labelsep=0.5em, itemsep=0.45em]
		\item[\textbf{第一步（代入乘积关系）：}]
			将 $a=4$，$x=1$ 代入 $xy = a^2 - a(x+y)$ 得
			\[
				\begin{aligned}
					1\cdot y &= 16 - 4(1+y) \\ &= 12 - 4y.
			\end{aligned}
			\]
			\par\textbf{理由：} 直接套用已知乘积关系。
		\item[\textbf{第二步（解一元方程）：}]
			移项合并得 $5y=12$，解得
			\[
				y = \frac{12}{5}.
			\]
			\par\textbf{理由：} 将方程整理为标准形式并求解。
		\item[\textbf{第三步（线段求和）：}]
			由 $EF = BE + DF$ 得
			\[
				\begin{aligned}
					EF &= 1 + \frac{12}{5} \\ &= \frac{17}{5}.
			\end{aligned}
			\]
			\par\textbf{理由：} 直接套用线段和结论。
	\end{description}
	\textbf{\textcolor{CExample}{关键结论：}}
	\textbf{$DF = \dfrac{12}{5}$，$EF = \dfrac{17}{5}$。}

	\medskip
	\textbf{\textcolor{CExample}{例 2（稍变形）}}
	\textbf{\textcolor{CExample}{题目：}}
	正方形 $ABCD$ 边长为 $5$，$E$ 在 $BC$ 上，$BE=2$，$\angle EAF=45^\circ$，求 $DF$、$EF$ 及 $\triangle CEF$ 的面积。
	\textbf{\textcolor{CExample}{解题步骤：}}
	\begin{description}[style=nextline, leftmargin=2.8em, labelsep=0.5em, itemsep=0.45em]
		\item[\textbf{第一步（代入求$y$）：}]
			将 $a=5$，$x=2$ 代入乘积关系得
			\[
				\begin{aligned}
					2y &= 25 - 5(2+y) \\ &= 15 - 5y.
			\end{aligned}
			\]
			\par\textbf{理由：} 代入已知量。
		\item[\textbf{第二步（解方程得$y$，再求$EF$）：}]
			移项得 $7y=15$，所以 $y=\dfrac{15}{7}$，进而
			\[
				\begin{aligned}
					EF &= 2 + \frac{15}{7} \\ &= \frac{29}{7}.
			\end{aligned}
			\]
			\par\textbf{理由：} 利用 $EF = BE + DF$。
		\item[\textbf{第三步（计算面积）：}]
			$EC = 5-2$，即 $EC=3$；$FC = 5-\dfrac{15}{7}$，即 $FC = \dfrac{20}{7}$，故
			\[
				\begin{aligned}
					S_{\triangle CEF} &= \frac12 \cdot 3 \cdot \frac{20}{7} \\ &= \frac{30}{7}.
			\end{aligned}
			\]
			\par\textbf{理由：} 直接应用直角三角形面积公式。
	\end{description}
	\textbf{\textcolor{CExample}{关键结论：}}
	\textbf{$DF = \dfrac{15}{7}$，$EF = \dfrac{29}{7}$，$S_{\triangle CEF} = \dfrac{30}{7}$。}

	\medskip
	\textbf{\textcolor{CExample}{例 3（综合应用）}}
	\textbf{\textcolor{CExample}{题目：}}
	正方形 $ABCD$ 边长为 $6$，$\angle EAF = 45^\circ$，$BE+DF=5$，求 $BE$ 和 $DF$ 的长。
	\textbf{\textcolor{CExample}{解题步骤：}}
	\begin{description}[style=nextline, leftmargin=2.8em, labelsep=0.5em, itemsep=0.45em]
		\item[\textbf{第一步（确定$EF$）：}]
			由 $EF = BE+DF$，得 $EF = 5$。
			\par\textbf{理由：} 直接应用线段和结论。
		\item[\textbf{第二步（求乘积$xy$）：}]
			代入 $xy = a^2 - a(x+y)$ 得
			\[
				\begin{aligned}
					xy &= 36 - 6 \times 5 \\ &= 6.
			\end{aligned}
			\]
			\par\textbf{理由：} 利用已知的乘积关系。
		\item[\textbf{第三步（逆用韦达定理）：}]
			$BE$ 和 $DF$ 是方程 $t^2 - 5t + 6 = 0$ 的两根，解得
			\[
				t_1 = 2,\quad t_2 = 3.
			\]
			\par\textbf{理由：} 已知和与积，可利用二次方程求解。
		\item[\textbf{第四步（得出结论）：}]
			因此 $BE=2$，$DF=3$，或 $BE=3$，$DF=2$。
			\par\textbf{理由：} 两组解都满足范围 $[0,6]$，均符合题意。
	\end{description}
	\textbf{\textcolor{CExample}{关键结论：}}
	\textbf{$BE=2$，$DF=3$ 或 $BE=3$，$DF=2$。}
\end{examplebox}
